\section{Reproducible evidence}
\subsection{Protocol and independent checks}
Results are regenerated by \texttt{python/experiments.py}, using local random
generators with seeds 11, 12 and 13. The code records Python, NumPy and SciPy
versions, the VICE revision, and SHA-256 hashes of its three imported modules.
SciPy provides numerical linear algebra and Qhull-backed triangulation
\cite{scipy2020,scipyDelaunay}. No inherited hard-coded benchmark numbers are
used. Default geometry uses centroids, $r=3$, a maximum edge-matrix condition
number of $10^8$, and explicit size guards; experiments use $M=15$ unless stated.

Exact affine tests cover $d=2,3,4$. Other tests cover translation, physical
chain-rule scaling, row permutations, tensor packing, plateau behavior, hull
flags, resource limits and an independent Gaussian sum. A Monte Carlo test
checks \cref{eq:covariance} using 150000 noise draws. These verify implementation
properties; they do not establish calibrated confidence.

\subsection{First-order accuracy: refinement and noise}
For dimensionless $\vect x\in[-1,1]^2$, use
\begin{equation}
 f(\vect x)=\tfrac12\vect x^{\mathsf T}\matr A\vect x+
             \vect b^{\mathsf T}\vect x,\qquad
 \matr A=\begin{bmatrix}1.2&0.2\\0.2&2.2\end{bmatrix},\quad
 \vect b=\begin{bmatrix}-0.4\\0.4\end{bmatrix}.
 \label{eq:quadratic}
\end{equation}
The analytic gradient $\matr A\vect x+\vect b$ is evaluated at retained
representatives. Each realization samples $N$ independent uniform points,
then adds Gaussian noise of absolute standard deviation $0$, $0.001$ or
$0.01$. Table statistics average the seed-wise statistics. The comparator
uses 24 nearest samples, a centered radius-scaled quadratic design matrix,
and ordinary least squares through the same observed values.

\begin{table}[htbp]
\centering
\caption{Mean seed-wise gradient error statistics. Noisy columns use absolute
noise standard deviation 0.01.}
\label{tab:accuracy}
\begin{tabular}{rrrrr}
\toprule
$N$ & Clean median & Clean 95th pct. & Noisy median & Quadratic noisy median\\
\midrule80 & 0.116 & 1.574 & 0.163 & 0.014 \\
160 & 0.072 & 0.803 & 0.164 & 0.018 \\
320 & 0.051 & 0.323 & 0.193 & 0.023 \\
640 & 0.035 & 0.143 & 0.260 & 0.031 \\
\bottomrule

\end{tabular}
\end{table}

Clean gradients improve with refinement, consistent with curvature bias.
At fixed noise, smaller cells amplify value error. The quadratic fit benefits
from additional observations and reproduces the exact quadratic without
noise to numerical precision. This comparator is especially favorable on
a quadratic; its superiority is not established for every objective.
Large tail errors also show why median accuracy alone can hide boundary
and slender-cell problems.

\begin{figure}[htbp]
\centering\includegraphics[width=\textwidth]{figures/generated/accuracy.pdf}
\caption{Left: means of seed-wise medians; bars show standard deviation over
three seeds, not confidence intervals. Right: two local estimators on the
same noisy samples. Refinement competes with noise amplification.}
\label{fig:accuracy}
\end{figure}

\subsection{Second-order recursion is a separate approximation}
For $N=160$, run the full two-stage scalar reconstruction. Restrict reported
medians to representatives with $\max_j|x_j|<0.8$, so the check is not solely
a hull artifact. The exact Hessian is $\matr A$ and Laplacian is $3.4$.
\begin{table}[htbp]
\centering
\caption{Mean seed-wise interior medians for recursive curvature. Matrix
error and symmetry defect use Frobenius norms.}
\label{tab:hessian}
\begin{tabular}{rrrr}
\toprule Noise std & Hessian error & Symmetry defect & Laplacian error\\
\midrule0.0 & 1.567 & 1.035 & 0.968 \\
0.001 & 1.613 & 1.044 & 1.015 \\
\bottomrule

\end{tabular}
\end{table}
The sizable clean error shows that first-order reproduction does not make
scalar recursion polynomial preserving. Differentiating an exact linear
vector gradient separately recovers $\matr A$ to numerical precision.
That test isolates the second-stage solve from first-stage support assignment.
Symmetrization cannot substitute for curvature validation.

\begin{figure}[htbp]
\centering\includegraphics[width=\textwidth]{figures/generated/hessian.pdf}
\caption{Clean quadratic, seed 11: recovered trace varies despite a constant
analytic Laplacian, and mixed partials disagree. Interior restriction does
not eliminate this discrepancy. Axes show the central 98\% ranges;
annotations count omitted tail cells.}
\label{fig:hessian}
\end{figure}

\subsection{Synthetic contamination: labels must follow the computation}
Sample $N=240$ points uniformly in $[-6,6]^2$ and evaluate
\begin{equation}
 f(x_1,x_2)=(x_1^2+x_2-11)^2+(x_1+x_2^2-7)^2.
\end{equation}
Add Gaussian noise of standard deviation $0.01\operatorname{std}(f)$.
Select 19 samples without replacement and add signed jumps
$1.4\operatorname{std}(f)(1+0.5U)$, with $U$ uniform on $[0,1]$.
The sample contamination fraction is about $7.9\%$.

A first gradient depends on all source vertices. Its target is therefore
``at least one source vertex was perturbed.'' Cell prevalence averages about
$21\%$. This denotes direct exposure, not a certified bad derivative.
Nearest-sample labels from the old demos can classify a cell as clean despite
a perturbed source vertex. They remain in saved arrays for audit, but are
not the reported target.

\begin{figure}[htbp]
\centering\includegraphics[width=\textwidth]{figures/generated/contamination.pdf}
\caption{Input and source-simplex labels live on different supports. The
descriptor is shown without interpreting it as probability of correctness.
Seed 11.}
\label{fig:contamination}
\end{figure}

Rank anomalies by $-S$ with the historical $\lambda=1$ flag threshold.
The comparator predicts each original sample from a quadratic fit of its
24 neighbors, excluding itself, and uses the maximum absolute residual among
each cell's source vertices. Both rankings use identical targets. At the
operating point, both flag exactly as many cells, providing a matched budget.
Neither ranking uses injected labels for tuning.

\begin{table}[htbp]
\centering
\caption{Mean synthetic screening metrics over three seeds. AP is average
precision from attained ranks, not trapezoidal PR area.}
\label{tab:screening}
\begin{tabular}{lrrrrr}
\toprule Method & Flag fraction & Precision & Recall & ROC AUC & AP\\
\midruleIso coverage & 0.914 & 0.177 & 0.769 & 0.372 & 0.176 \\
Local quadratic residual & 0.914 & 0.229 & 1.000 & 0.999 & 0.997 \\
\bottomrule

\end{tabular}
\end{table}
Low coverage is a poor anomaly ranking here: mean AUC is below $0.5$,
and AP below mean cell prevalence. Across seeds, iso AUC standard deviation
is about $0.178$; this small sample does not justify a universal negative
theorem. It does refute a claim that these demonstrations establish reliable
outlier detection. The strong residual comparator applies to large injected
jumps, not necessarily subtle measurement errors.

\begin{figure}[htbp]
\centering\includegraphics[width=\textwidth]{figures/generated/screening.pdf}
\caption{Exact rank sweeps for seed 11 with tied ranks and ROC endpoints.
The precision-recall plot retains the cell prevalence reference, relevant for
imbalanced targets \cite{saito2015}. Simplex centroids are not independent
validation observations.}
\label{fig:screening}
\end{figure}

Sensitivity results vary $M\in\{8,15,30\}$ and bandwidth factors
$\{0.5,1,2\}$ on seed 11, check quantile levels, and vary
$\lambda\in\{-1,0,1,2\}$. None of these level/bandwidth settings repairs the
low-score ranking on this realization. Raising $\lambda$ increases flag counts,
as the formula requires. This finite sweep is not an optimized detector search.

\subsection{Search proposals: measure actual gain}
For $N=240$ samples of \cref{eq:quadratic}, evaluate descent trials on the
analytic objective. Report proposals inside the input convex hull, and
separately those in the highest coverage quartile.
\begin{table}[htbp]
\centering
\caption{Mean seed-wise fractions of strict improvement. The last column
averages seed-wise Spearman correlations of coverage with gradient error.}
\label{tab:guidance}
\begin{tabular}{rrrr}
\toprule Noise std & All inside & High coverage inside & Coverage--error correlation\\
\midrule0.0 & 1.000 & 1.000 & 0.049 \\
0.01 & 0.991 & 0.997 & 0.050 \\
\bottomrule

\end{tabular}
\end{table}
These short trials usually improve this smooth objective. Weak coverage--error
correlation does not calibrate accuracy. Both the high-coverage subset and
overall population already have high success fractions; useful incremental
benefit from gating is not established. Synthetic true-objective evaluation
must become an actual measurement in a real loop. This tests one-step trials,
not trajectory convergence.

\begin{figure}[htbp]
\centering\includegraphics[width=\textwidth]{figures/generated/guidance.pdf}
\caption{Seed 11, noise std 0.01. Left: a subset of descent proposals. Right:
coverage versus analytic gradient error. A geometric descriptor needs an
independent accuracy target.}
\label{fig:guidance}
\end{figure}

\subsection{Dimension, cost and support limits}
Small one-stage quadratic checks at $N=50$ in $d=2,3,4$ produce 84, 224 and
617 representatives, with median gradient errors approximately 0.137, 0.599
and 1.518. These are feasibility checks, not scalability benchmarks.
Increasing dimension at fixed $N$ also makes support sparser.

For $T$ current simplices and $C=d^{k-1}$ value components, solves cost about
$O(Td^3+TCd^2)$; the next support has up to $T$ points. The edge-level stage
allocates an $E\times M$ interpolation matrix per component. Scoring $Q$
queries against $P_c$ cuts costs $O(QP_cd)$ per component; chunks bound
temporary memory without changing this work. Both support and component
counts can grow rapidly under recursion.

The research limits are 100000 simplices and two million potential edge-level
pairs per component. Guards run after triangulation or edge enumeration, so
they do not guarantee Qhull fits a fixed memory budget. Higher-dimensional
industrial maps require another neighborhood strategy or resource planning.
Rank-deficient support, duplicate representatives, nonfinite values and
overstrict edge filtering are explicit failures. Qhull omissions are recorded
and warned about rather than silently treated as fully supported samples.
